
\begin{frame}{Definition 3: Calibration and Predictive Parity (Sufficiency)}

    \begin{block}{Calibration within groups}
        \[
            \Pr(Y = 1 \mid R = r,\, A = \GrpA) \;=\; \Pr(Y = 1 \mid R = r,\, A = \GrpB) \;=\; r
        \]
        A score of $0.7$ means a 70\% chance for everyone, whatever their group.
        \textbf{Predictive parity} is the thresholded version: equal PPV across groups.
    \end{block}

    \vspace{0.8em}

    \begin{itemize}
        \setlength\itemsep{0.45em}
        \item This is the criterion the COMPAS vendor pointed to, and it is a genuinely important
              one: without it, a human reading the score is being misled about what it means.
        \item Well-trained probabilistic models are often \emph{approximately} calibrated
              within groups by default --- which is precisely why the conflict below is unavoidable
              rather than exotic.
        \item Check it with a \textbf{reliability diagram per group}, not with a single number.
    \end{itemize}
\end{frame}
